\section{Point Estimation Methods}

In the previous section we defined an \emph{estimator} as a function of the random sample $X_1, X_2, \ldots, X_n$ used to infer the value of an unknown population parameter $\theta \in \Theta$. While basic theory allows us to evaluate whether a proposed estimator possesses desirable properties (such as unbiasedness or low variance), in practice we need systematic mathematical procedures to \emph{construct} or derive such estimators from a statistical model.

In this section we will study the two most fundamental point estimation methods in modern statistical theory: the \emph{Method of Moments} (MoM) and the \emph{Maximum Likelihood Estimation} (MLE) method. We will also analyze their roles in model design and training in data science and machine learning.

\subsection{Quality Criteria for a Point Estimator}

Before studying how to \emph{construct} estimators, we formalize the criteria used to \emph{evaluate} them: unbiasedness, efficiency, and consistency.

\begin{definicion}[Point estimator]
Let $X_1, \ldots, X_n$ be a random sample from a population with unknown parameter $\theta \in \Theta$. A \emph{point estimator} of $\theta$ is any statistic $\hat{\theta} = T(X_1, \ldots, X_n)$, i.e., a function of the sample that does not depend on $\theta$. Being a function of random variables, $\hat{\theta}$ is itself a random variable with its own sampling distribution.
\end{definicion}

\begin{definicion}[Bias]
 \label{eq:6.1.1}
The \emph{bias} of an estimator $\hat{\theta}$ with respect to the parameter $\theta$ is defined as
\begin{align}
 \text{Bias}(\hat{\theta}) = E(\hat{\theta}) - \theta.
\end{align}
If $\text{Bias}(\hat{\theta}) = 0$ for all $\theta \in \Theta$, we say $\hat{\theta}$ is \emph{unbiased}.
\end{definicion}

\begin{definicion}[Mean Squared Error]
 \label{eq:6.1.2}
The \emph{mean squared error} (MSE) of an estimator is
\begin{align}
 \text{MSE}(\hat{\theta}) = E\left[(\hat{\theta}-\theta)^2\right].
\end{align}
\end{definicion}

\begin{teorema}[Bias-variance decomposition]
 \label{eq:6.1.3}
For any estimator $\hat{\theta}$,
\begin{align}
 \text{MSE}(\hat{\theta}) = \text{Var}(\hat{\theta}) + \left[\text{Bias}(\hat{\theta})\right]^2.
\end{align}
\end{teorema}

\begin{proof}
Adding and subtracting $E(\hat{\theta})$ inside the square:
\begin{align}
 E\left[(\hat{\theta}-\theta)^2\right] &= E\left[\left((\hat{\theta}-E(\hat{\theta})) + (E(\hat{\theta})-\theta)\right)^2\right] \nonumber \\
 &= E\left[(\hat{\theta}-E(\hat{\theta}))^2\right] + 2(E(\hat{\theta})-\theta)\underbrace{E\left[\hat{\theta}-E(\hat{\theta})\right]}_{=0} + (E(\hat{\theta})-\theta)^2 \nonumber \\
 &= \text{Var}(\hat{\theta}) + \left[\text{Bias}(\hat{\theta})\right]^2. \qedhere
\end{align}
\end{proof}

\begin{observacion}[The bias-variance trade-off]
This decomposition explains why a slightly biased estimator can sometimes be preferable to an unbiased one: if the introduced bias is small but substantially reduces variance, the total MSE can decrease. This \emph{bias-variance trade-off} is central to regularization in machine learning models (Ridge, Lasso), where a controlled bias is accepted in exchange for reduced variance and, therefore, better generalization error.
\end{observacion}

\begin{definicion}[Relative efficiency]
 \label{eq:6.1.4}
If $\hat{\theta}_1$ and $\hat{\theta}_2$ are two unbiased estimators of $\theta$, the \emph{relative efficiency} of $\hat{\theta}_1$ with respect to $\hat{\theta}_2$ is
\begin{align}
 \text{Eff}(\hat{\theta}_1, \hat{\theta}_2) = \frac{\text{Var}(\hat{\theta}_2)}{\text{Var}(\hat{\theta}_1)}.
\end{align}
If $\text{Eff}(\hat{\theta}_1,\hat{\theta}_2) > 1$, $\hat{\theta}_1$ is more efficient (has lower variance).
\end{definicion}

\begin{teorema}[Cramér-Rao Lower Bound]
 \label{eq:6.1.5}
Let $\hat{\theta}$ be an unbiased estimator of $\theta$ based on an iid sample of size $n$, and let $I(\theta) = -E\left[\dfrac{\partial^2}{\partial\theta^2}\ln f(X\mid\theta)\right]$ be the Fisher Information of a single observation. Under regularity conditions,
\begin{align}
 \text{Var}(\hat{\theta}) \geq \frac{1}{n\,I(\theta)}.
\end{align}
An unbiased estimator that attains this bound exactly is called a \emph{Uniformly Minimum Variance Unbiased Estimator} (UMVUE) and is considered \emph{efficient} in the absolute sense.
\end{teorema}

\begin{definicion}[Consistency]
 \label{eq:6.1.6}
A sequence of estimators $\hat{\theta}_n$ (indexed by sample size $n$) is \emph{consistent} for $\theta$ if it converges in probability to the true parameter:
\begin{align}
 \hat{\theta}_n \xrightarrow{P} \theta \quad \text{as } n \to \infty,
\end{align}
i.e., for every $\varepsilon > 0$, $\lim_{n\to\infty} P(|\hat{\theta}_n - \theta| \geq \varepsilon) = 0$.
\end{definicion}

\begin{observacion}[A sufficient condition for consistency]
By Chebyshev's inequality, $P(|\hat{\theta}_n-\theta|\ge\varepsilon) \le \text{MSE}(\hat{\theta}_n)/\varepsilon^2$. Therefore, if $\text{MSE}(\hat{\theta}_n) \to 0$ as $n\to\infty$ (equivalently, if both the bias and variance tend to zero), then $\hat{\theta}_n$ is consistent. This is the most common way to prove consistency in practice.
\end{observacion}

\begin{ejemplo}
 \label{exmp:6.1.1}
For an iid sample $X_1,\ldots,X_n$ from a population with mean $\mu$ and variance $\sigma^2$, compare the MSE of two estimators of $\mu$: the sample mean $\bar{X}$ and the naive single-observation estimator $T_1 = X_1$.
\end{ejemplo}

\begin{solucion}
Both are unbiased: $E(\bar{X})=\mu$ and $E(T_1)=\mu$, so $\text{MSE}=\text{Var}$ in both cases.
\begin{align}
 \text{Var}(\bar{X}) = \frac{\sigma^2}{n}, \qquad \text{Var}(T_1) = \sigma^2.
\end{align}
The relative efficiency is $\text{Eff}(\bar{X}, T_1) = \sigma^2/(\sigma^2/n) = n$: the sample mean is $n$ times more efficient than using a single observation, and it is consistent ($\text{Var}(\bar{X})\to 0$), while $T_1$ is not (its variance does not depend on $n$).
\end{solucion}

\subsection{Method of Moments (MoM)}

The \emph{Method of Moments}, proposed by Pafnuty Chebyshev and Karl Pearson in the late 19th century, is one of the most intuitive procedures in statistical inference. It is based on equating the numerical characteristics of the sample (the sample moments) with the theoretical or expected moments of the population.

\begin{definicion}[Population and Sample Moments]
	Let $X$ be a random variable with density function $f(x \mid \theta)$. The $k$-th \emph{population moment about the origin}, denoted by $\mu_k(\theta)$, is defined as:
	\begin{align}
		\mu_k(\theta) = E(X^k) = \int_{-\infty}^{\infty} x^k f(x \mid \theta) \, dx \quad \text{(continuous case)}.
	\end{align}

	Given a random sample $X_1, X_2, \ldots, X_n$, the $k$-th \emph{sample moment}, denoted by $m_k$, is the average of the $k$-th powers of the data:
	\begin{align}
		m_k = \frac{1}{n} \sum_{i=1}^n X_i^k.
	\end{align}
	In particular, the first sample moment is the arithmetic mean, $m_1 = \bar{X}$, and the second sample moment is $m_2 = \frac{1}{n}\sum X_i^2$.
\end{definicion}

\begin{teorema}[Method of Moments Procedure]
	Suppose the population distribution depends on a vector of $r$ unknown parameters, $\theta = (\theta_1, \theta_2, \ldots, \theta_r)$. The estimation procedure by the method of moments consists of equating the first $r$ population moments to the respective $r$ sample moments:
	\begin{align}
		\mu_1(\theta_1, \ldots, \theta_r) &= m_1 = \bar{X} \nonumber \\
		\mu_2(\theta_1, \ldots, \theta_r) &= m_2 = \frac{1}{n}\sum_{i=1}^n X_i^2 \nonumber \\
		&\vdots \nonumber \\
		\mu_r(\theta_1, \ldots, \theta_r) &= m_r = \frac{1}{n}\sum_{i=1}^n X_i^r.
		\label{eq:mom_sistema}
	\end{align}

	The simultaneous solution of this algebraic system of $r$ equations in $r$ unknowns produces the vector of method of moments estimators, denoted by $\hat{\theta}_{MoM} = (\hat{\theta}_{1, MoM}, \ldots, \hat{\theta}_{r, MoM})$.
\end{teorema}

\begin{ejemplo}
	\textbf{Method of Moments for the Gamma Distribution.}
	A random variable $X \sim \text{Gamma}(\alpha, \beta)$ has density $f(x) = \frac{1}{\beta^\alpha \Gamma(\alpha)} x^{\alpha-1} e^{-x/\beta}$ for $x > 0$. As known from the continuous random variables unit, the theoretical mean and variance of the gamma distribution are:
	\begin{align}
		E(X) &= \mu_1 = \alpha\beta, \\
		\text{Var}(X) &= \alpha\beta^2.
	\end{align}

	Since $\text{Var}(X) = E(X^2) - [E(X)]^2 = \mu_2 - \mu_1^2$, we can express the first two population moments as:
	\begin{align}
		\mu_1 &= \alpha\beta, \\
		\mu_2 &= \text{Var}(X) + \mu_1^2 = \alpha\beta^2 + \alpha^2\beta^2 = \alpha\beta^2(1 + \alpha).
	\end{align}

	Equating with the sample moments $m_1 = \bar{X}$, we can equivalently work by equating the sample variance $S^2$ with the population variance:
	\begin{align}
		\alpha\beta &= \bar{X}, \\
		\alpha\beta^2 &= S^2 \quad \text{(where } S^2 = m_2 - m_1^2 \text{)}.
	\end{align}

	Dividing the second equation by the first, we obtain the estimator for the scale parameter $\beta$:
	\begin{align}
		\frac{\alpha\beta^2}{\alpha\beta} = \frac{S^2}{\bar{X}} \implies \hat{\beta}_{MoM} = \frac{S^2}{\bar{X}}.
	\end{align}

	Substituting $\hat{\beta}_{MoM}$ into the first equation yields the estimator for the shape parameter $\alpha$:
	\begin{align}
		\alpha \left( \frac{S^2}{\bar{X}} \right) = \bar{X} \implies \hat{\alpha}_{MoM} = \frac{\bar{X}^2}{S^2}.
	\end{align}

	This is a practically useful result because the maximum likelihood system of equations for the Gamma distribution does not have a closed-form analytical solution and requires iterative numerical methods (such as Newton-Raphson), whereas the method of moments provides immediate closed-form formulas.
\end{ejemplo}

\begin{ejemplo}
	\textbf{A delicate case: when the first moment does not identify the parameter.}
	Let $X \sim U(-\theta, \theta)$ with unknown $\theta > 0$. By the symmetry of the distribution, the first population moment is identically zero for any value of $\theta$:
	\begin{align}
		\mu_1 = E(X) = 0, \quad \text{for all } \theta > 0.
	\end{align}
	Equating $\mu_1 = m_1 = \bar{X}$ does not allow us to solve for $\theta$: the first moment simply carries no information about the parameter. We must resort to the first \emph{non-trivial} moment, in this case the second moment about the origin. Since $\text{Var}(X) = \theta^2/3$ and $\mu_1 = 0$:
	\begin{align}
		\mu_2 = \text{Var}(X) + \mu_1^2 = \frac{\theta^2}{3}.
	\end{align}
	Equating with the second sample moment $m_2 = \frac{1}{n}\sum X_i^2$:
	\begin{align}
		\frac{\theta^2}{3} = m_2 \implies \hat{\theta}_{MoM} = \sqrt{3 m_2} = \sqrt{\frac{3}{n}\sum_{i=1}^n X_i^2}.
	\end{align}
\end{ejemplo}

\begin{observacion}[Properties and limitations of the Method of Moments]
	MoM has clear advantages and disadvantages relative to MLE:
	\begin{itemize}
		\item \textbf{Consistency:} by the Law of Large Numbers, $m_k \xrightarrow{P} \mu_k(\theta)$ for each $k$; under regularity conditions (continuity of the function relating moments to parameters), this implies $\hat{\theta}_{MoM} \xrightarrow{P} \theta$.
		\item \textbf{Computational simplicity:} it often yields closed-form formulas (as in the Gamma example above) even when MLE requires iterative numerical optimization.
		\item \textbf{Lower efficiency:} in general, $\hat{\theta}_{MoM}$ does not attain the Cramer-Rao Bound and has higher asymptotic variance than $\hat{\theta}_{MLE}$; for this reason, MoM is frequently used only as a starting value for iterative MLE algorithms (see data science applications below).
		\item \textbf{Estimates outside the parameter space:} unlike MLE (which by construction always yields an admissible value), MoM can occasionally produce estimates outside $\Theta$ (for example, a negative estimated variance in models with nonlinear restrictions between moments and parameters).
	\end{itemize}
\end{observacion}

\subsection{Maximum Likelihood Estimation (MLE)}

The principle of maximum likelihood, formally developed by Ronald A. Fisher in the early 20th century, holds that, given a set of sample observations, we should choose as our estimate that value of the parameter $\theta$ which makes the observed sample as probable as possible.

\begin{definicion}[Likelihood Function]
	Let $X_1, X_2, \ldots, X_n$ be a random sample of size $n$ from a population with probability density function (or probability mass function if discrete) $f(x \mid \theta)$, where $\theta$ is an unknown parameter in the parameter space $\Theta$.

	Since the observations are independent and identically distributed (\emph{iid}), the joint density of the sample can be expressed as the product of the marginal densities. Evaluated at the observed values $x_1, x_2, \ldots, x_n$, this function viewed as a function of the parameter $\theta$ is called the \emph{likelihood function} and is denoted by $L(\theta \mid x_1, \ldots, x_n)$ or simply $L(\theta)$:
	\begin{align}
		L(\theta) = f(x_1, x_2, \ldots, x_n \mid \theta) = \prod_{i=1}^n f(x_i \mid \theta).
		\label{eq:verosimilitud_def}
	\end{align}
\end{definicion}

\begin{observacion}
	It is fundamental to distinguish between the joint density and the likelihood function. In the joint density function $f(x_1, \ldots, x_n \mid \theta)$, the parameter $\theta$ is considered fixed and known, while the variables $X_i$ are variable. Conversely, in the likelihood function $L(\theta)$, the observed data $x_1, \ldots, x_n$ are fixed and known, and $\theta$ is the free variable.
\end{observacion}

\begin{definicion}[Maximum Likelihood Estimator]
	The \emph{maximum likelihood estimator} (MLE) of the parameter $\theta$, denoted by $\hat{\theta}_{MLE}$, is the value in the parameter space $\Theta$ that maximizes the likelihood function:
	\begin{align}
		L(\hat{\theta}_{MLE}) \geq L(\theta) \quad \text{for all } \theta \in \Theta,
	\end{align}
	or equivalently:
	\begin{align}
		\hat{\theta}_{MLE} = \arg\max_{\theta \in \Theta} L(\theta).
	\end{align}
\end{definicion}

Because the natural logarithm function $\ln(\cdot)$ is strictly increasing on its positive domain, maximizing $L(\theta)$ is mathematically equivalent to maximizing its logarithm. We commonly work with the \emph{log-likelihood function}, denoted by $\ell(\theta)$:
\begin{align}
	\ell(\theta) = \ln L(\theta) = \ln \left( \prod_{i=1}^n f(x_i \mid \theta) \right) = \sum_{i=1}^n \ln f(x_i \mid \theta).
	\label{eq:log_verosimilitud}
\end{align}

The logarithmic transformation converts products into sums, which greatly simplifies analytical differentiation and prevents numerical underflow in computational implementations.

\begin{teorema}[First-Order Condition for MLE]
	If the log-likelihood function $\ell(\theta)$ is differentiable with respect to $\theta$ and the maximum is attained in the interior of the parameter space $\Theta$, the maximum likelihood estimator $\hat{\theta}_{MLE}$ satisfies the score equation:
	\begin{align}
		\frac{d}{d\theta} \ell(\theta) = \sum_{i=1}^n \frac{d}{d\theta} \ln f(x_i \mid \theta) = 0.
		\label{eq:mle_score}
	\end{align}
	To verify that the solution corresponds to a global maximum and not a minimum or inflection point, one must confirm that the second derivative is strictly negative at that point:
	\begin{align}
		\left. \frac{d^2}{d\theta^2} \ell(\theta) \right|_{\theta = \hat{\theta}_{MLE}} < 0.
	\end{align}
\end{teorema}

\subsubsection{The Score Function and Asymptotic Normality}

The quantity set to zero in the first-order condition has its own name and importance: it is the \emph{score function}.

\begin{definicion}[Score function]
	For a single observation $X$ with density $f(x\mid\theta)$, the \emph{score function} is
	\begin{align}
		U(\theta; X) = \frac{\partial}{\partial\theta}\ln f(X\mid\theta).
	\end{align}
	For an iid sample $X_1,\ldots,X_n$, the \emph{total score} is $U_n(\theta) = \sum_{i=1}^n U(\theta; X_i) = \dfrac{d}{d\theta}\ell(\theta)$. The maximum likelihood estimator is, by definition, the solution of the \emph{score equation} $U_n(\hat\theta_{MLE}) = 0$.
\end{definicion}

\begin{propiedad}[Properties of the score]
	Under regularity conditions permitting the interchange of differentiation and integration:
	\begin{enumerate}
		\item \textbf{Zero mean:} $E_\theta[U(\theta;X)] = 0$ at the true value of $\theta$.
		\item \textbf{Information identity:} $\text{Var}_\theta[U(\theta;X)] = I(\theta)$, i.e., the variance of a single observation's score is exactly the Fisher Information.
	\end{enumerate}
\end{propiedad}

\begin{proof}[Proof sketch of (1)]
	Since $\int f(x\mid\theta)\,dx=1$ for all $\theta$, differentiating both sides with respect to $\theta$ and interchanging derivative and integral:
	\begin{align}
		0 = \frac{\partial}{\partial\theta}\int f(x\mid\theta)\,dx = \int \frac{\partial f(x\mid\theta)/\partial\theta}{f(x\mid\theta)} f(x\mid\theta)\,dx = E_\theta[U(\theta;X)]. \qedhere
	\end{align}
\end{proof}

\begin{teorema}[Asymptotic normality of the MLE]
	 \label{eq:6.3.1}
	Under regularity conditions, as $n\to\infty$,
	\begin{align}
		\sqrt{n}\left(\hat\theta_{MLE} - \theta\right) \xrightarrow{d} N\!\left(0,\, \frac{1}{I(\theta)}\right),
	\end{align}
	i.e., for large $n$, $\hat\theta_{MLE} \overset{\text{approx.}}{\sim} N\!\left(\theta,\, \dfrac{1}{nI(\theta)}\right)$.
\end{teorema}

\begin{observacion}[Connection to the Cramer-Rao Bound]
	This result formalizes the \emph{asymptotic efficiency} of the MLE mentioned in Section 06.01: the asymptotic variance of the MLE, $1/(nI(\theta))$, coincides exactly with the Cramer-Rao Lower Bound. Thus, even though the MLE may be neither unbiased nor optimal for finite $n$ (as we saw with $\hat\sigma^2_{MLE}$), it becomes asymptotically optimal: no consistent estimator can do better in the limit.
\end{observacion}

\begin{ejemplo}
	\textbf{MLE for the Bernoulli Distribution.}
	Suppose $X_1, X_2, \ldots, X_n$ is a random sample of Bernoulli random variables with parameter $p$, where $P(X_i = 1) = p$ and $P(X_i = 0) = 1-p$. The probability mass function for a single observation is:
	\begin{align}
		f(x_i \mid p) = p^{x_i}(1-p)^{1-x_i}, \quad x_i \in \{0, 1\}.
	\end{align}

	The likelihood function for the sample is:
	\begin{align}
		L(p) = \prod_{i=1}^n p^{x_i}(1-p)^{1-x_i} = p^{\sum x_i}(1-p)^{n - \sum x_i}.
	\end{align}

	Taking the natural logarithm, we obtain the log-likelihood:
	\begin{align}
		\ell(p) = \left(\sum_{i=1}^n x_i\right) \ln p + \left(n - \sum_{i=1}^n x_i\right) \ln(1-p).
	\end{align}

	Differentiating with respect to the parameter $p$ and setting equal to zero:
	\begin{align}
		\frac{d}{dp} \ell(p) = \frac{\sum x_i}{p} - \frac{n - \sum x_i}{1-p} = 0.
	\end{align}

	Multiplying both sides by $p(1-p)$ and simplifying:
	\begin{align}
		(1-p)\sum_{i=1}^n x_i - p\left(n - \sum_{i=1}^n x_i\right) &= 0 \\
		\sum_{i=1}^n x_i - p\sum_{i=1}^n x_i - np + p\sum_{i=1}^n x_i &= 0 \\
		\sum_{i=1}^n x_i - np &= 0 \implies \hat{p}_{MLE} = \frac{1}{n} \sum_{i=1}^n X_i = \bar{X}.
	\end{align}

	Therefore, the maximum likelihood estimator for the probability of success $p$ in a Bernoulli (and Binomial) distribution is the sample proportion of successes, $\bar{X}$.
\end{ejemplo}

\begin{ejemplo}
	\textbf{MLE for the Parameters of the Normal Distribution.}
	Let $X_1, X_2, \ldots, X_n$ be a random sample from a normal population with unknown mean $\mu$ and variance $\sigma^2$. The joint density is:
	\begin{align}
		f(x_i \mid \mu, \sigma^2) = \frac{1}{\sqrt{2\pi\sigma^2}} \exp\left( -\frac{(x_i - \mu)^2}{2\sigma^2} \right).
	\end{align}

	The log-likelihood function, with parameter vector $\theta = (\mu, \sigma^2)$, is given by:
	\begin{align}
		\ell(\mu, \sigma^2) &= \sum_{i=1}^n \left[ -\frac{1}{2}\ln(2\pi) - \frac{1}{2}\ln(\sigma^2) - \frac{(x_i - \mu)^2}{2\sigma^2} \right] \\
		&= -\frac{n}{2}\ln(2\pi) - \frac{n}{2}\ln(\sigma^2) - \frac{1}{2\sigma^2}\sum_{i=1}^n (x_i - \mu)^2.
	\end{align}

	To find the partial estimates, we differentiate with respect to $\mu$ and $\sigma^2$:
	\begin{align}
		\frac{\partial}{\partial \mu} \ell(\mu, \sigma^2) = \frac{1}{\sigma^2}\sum_{i=1}^n (x_i - \mu) = 0 \implies \sum_{i=1}^n x_i - n\mu = 0 \implies \hat{\mu}_{MLE} = \bar{X}.
	\end{align}

	Now differentiating with respect to the variance $\sigma^2$ (treating $\sigma^2$ as the variable of differentiation):
	\begin{align}
		\frac{\partial}{\partial (\sigma^2)} \ell(\mu, \sigma^2) = -\frac{n}{2\sigma^2} + \frac{1}{2(\sigma^2)^2} \sum_{i=1}^n (x_i - \mu)^2 = 0.
	\end{align}

	Substituting $\hat{\mu}_{MLE} = \bar{X}$ and multiplying by $2(\sigma^2)^2$:
	\begin{align}
		-n\sigma^2 + \sum_{i=1}^n (x_i - \bar{X})^2 = 0 \implies \hat{\sigma}^2_{MLE} = \frac{1}{n}\sum_{i=1}^n (X_i - \bar{X})^2.
	\end{align}
\end{ejemplo}

\begin{observacion}
	Notice that the maximum likelihood estimator for the normal variance, $\hat{\sigma}^2_{MLE} = \frac{1}{n}\sum (X_i - \bar{X})^2$, divides by $n$ rather than $n-1$. As demonstrated in the sampling distributions chapter, this implies that $\hat{\sigma}^2_{MLE}$ is a slightly downward-biased estimator ($E(\hat{\sigma}^2_{MLE}) = \frac{n-1}{n}\sigma^2$). However, as $n \to \infty$, the bias vanishes rapidly and the estimator is fully consistent.
\end{observacion}

\begin{propiedad}[Properties of Maximum Likelihood Estimators]
	Maximum likelihood estimators enjoy excellent theoretical properties under regularity conditions:
	\begin{enumerate}
		\item \textbf{Functional Invariance:} If $\hat{\theta}$ is the MLE of $\theta$ and $g(\theta)$ is any one-to-one (or measurable) function of $\theta$, then the MLE of $g(\theta)$ is exactly $\widehat{g(\theta)} = g(\hat{\theta})$. For example, the MLE of the normal standard deviation is $\hat{\sigma}_{MLE} = \sqrt{\hat{\sigma}^2_{MLE}}$.
		\item \textbf{Consistency:} As the sample size grows, $\hat{\theta}_{MLE}$ converges in probability to the true value of the parameter: $\hat{\theta}_{MLE} \xrightarrow{p} \theta$.
		\item \textbf{Asymptotic Normality:} For large samples, the sampling distribution of $\hat{\theta}_{MLE}$ approaches a normal distribution centered at the true value $\theta$ with variance equal to the inverse of the Fisher Information.
		\item \textbf{Asymptotic Efficiency:} Among all consistent and asymptotically normal estimators, the MLE achieves the minimum theoretical variance possible as $n \to \infty$ (the Cramér-Rao lower bound).
	\end{enumerate}
\end{propiedad}

\subsection{Point Estimation Methods in Data Science}

In modern data science and machine learning, point estimation methods form the mathematical core of model training algorithms:

\begin{enumerate}
	\item \textbf{Computational Training via MLE:} Many supervised learning algorithms, such as logistic regression, Poisson regression for count data, and deep neural networks, are trained by directly maximizing a likelihood function. For example, in binary or multi-class classification, minimizing the \emph{cross-entropy loss} function is mathematically identical to computing the maximum likelihood estimator for a multinomial distribution.

	\item \textbf{Weight Initialization with Method of Moments:} When fitting complex statistical models or mixture distributions (e.g., Gaussian Mixture Models - GMM) using intensive numerical optimization methods like the Expectation-Maximization (EM) algorithm, estimates from the Method of Moments are frequently used as the initial starting values ($\theta^{(0)}$) to accelerate convergence and avoid poor local minima.

	\item \textbf{Generalized Method of Moments (GMM) Estimation:} In econometrics and causal data science, when the number of moment conditions exceeds the number of parameters to estimate, the Generalized Method of Moments is employed to optimally combine all available empirical information without requiring strict distributional specifications.
\end{enumerate}
